\section{Kesimpulan}
Program nomor 1 membaca citra, menghitung dua respons LoG pada
$\sigma=0{,}1$ dan $\sigma=0{,}2$ piksel, lalu menyimpan selisih
\texttt{im1 - im2} sebagai data float64 bertanda. Visualisasi dibuat setelah
pengurangan, sehingga penyesuaian tampilan tidak mengubah hasil numerik.
Sigma yang sangat kecil menyebabkan bias respons pada grid diskret;
hal tersebut perlu dibedakan dari perilaku ideal LoG kontinu.

Program nomor 2 memakai komponen terhubung dan \kode{regionprops}
untuk memisahkan 18 karakter pada citra contoh menjadi berkas individual.
Mask tiap region mempertahankan orientasi karakter dan menghindari
penyalinan karakter tetangga di dalam kotak pembatas.

Program nomor 3 memisahkan koin melalui ambang Otsu dan komponen
terhubung, mengukur luas serta diameter ekuivalen menggunakan
\kode{regionprops}, kemudian menghitung jumlah tiap kelompok ukuran.
Pada contoh sintetis, diperoleh 100 koin dalam empat kelompok, masing-masing
25 buah. Delapan pengujian lulus. Penerapan bergantung pada asumsi koin
terpisah, terlihat utuh, dan kontras terhadap latar seragam.

\clearpage
\begin{thebibliography}{9}
\raggedright
\small
\setstretch{1}
\setlength{\itemsep}{3pt}
\setlength{\parskip}{0pt}
\addcontentsline{toc}{section}{Daftar Pustaka}
\bibitem{soal-uts}
Peter Rhatodirdjo Angka,
\textit{UTS INF301 Pengolahan Citra Digital, Semester Gasal 2026/2027}.
Program Studi Informatika, Fakultas Teknik, Universitas Katolik Widya
Mandala Surabaya, 8 Oktober 2026. Dokumen soal yang diberikan.

\bibitem{scipy-log}
SciPy Developers,
\textit{scipy.ndimage.gaussian\_laplace}.
Dokumentasi resmi SciPy.
\href{https://docs.scipy.org/doc/scipy/reference/generated/scipy.ndimage.gaussian_laplace.html}{Tersedia pada laman resmi SciPy}.

\bibitem{skimage-regions}
scikit-image Developers,
\textit{skimage.measure.regionprops}.
Dokumentasi resmi scikit-image.
\href{https://scikit-image.org/docs/stable/api/skimage.measure.html#skimage.measure.regionprops}{Tersedia pada laman resmi scikit-image}.

\bibitem{skimage-otsu}
scikit-image Developers,
\textit{skimage.filters.threshold\_otsu}.
Dokumentasi resmi scikit-image.
\href{https://scikit-image.org/docs/stable/api/skimage.filters.html#skimage.filters.threshold_otsu}{Tersedia pada laman resmi scikit-image}.
\end{thebibliography}
